\section*{Общий блок · вопросы 1–6}
\begin{enumerate}
\setlength{\itemsep}{4mm}
\question{Какую роль выполняет CPU при работе программы?}
  {Постоянно хранит исходный код.}{Выполняет машинные инструкции.}
  {Заменяет оперативную память.}{Сохраняет файлы после выключения питания.}
\question{Где находится программа во время исполнения?}
  {Её инструкции загружены в RAM, откуда их получает CPU.}
  {Она существует только как текст исходного файла.}
  {Она целиком помещается в один регистр.}
  {Она исполняется непосредственно с SSD, минуя RAM.}
\question{Чем адрес ячейки памяти отличается от её значения?}
  {Адрес обозначает место, значение — данные в этом месте.}
  {Адрес всегда равен значению.}
  {Значение обозначает место, адрес — данные в нём.}
  {Адрес есть только у инструкций, но не у данных.}
\question{Что такое регистр процессора?}
  {Область хранения файлов на накопителе.}
  {Небольшая рабочая ячейка внутри CPU.}
  {Любая переменная в исходном коде.}
  {Ячейка видеопамяти.}
\question{Как расположены элементы массива \texttt{unsigned char a[4]} в памяти?}
  {Каждый элемент обязательно находится в отдельном регистре.}
  {Элементы следуют друг за другом.}
  {Адреса элементов выбираются случайно.}
  {Элементы всегда хранятся на диске.}
\question{Для чего программе нужен стек при вызове функции?}
  {Чтобы хранить адрес возврата и другие данные вызова.}
  {Чтобы сохранять файлы после выключения питания.}
  {Чтобы заменить код функции машинными инструкциями.}
  {Чтобы выводить символы на экран.}
\end{enumerate}
